\documentclass[tikz,border=6pt]{standalone}
\usepackage{xeCJK}
\usepackage{fontspec}
\setmainfont{Times New Roman}
\setCJKmainfont{Microsoft YaHei}
\usepackage{amsmath}
\usetikzlibrary{arrows.meta,positioning,fit,calc,shapes.multipart,backgrounds}

\begin{document}
\begin{tikzpicture}[
  font=\small,
  >=Latex,
  node distance=8mm and 10mm,
  box/.style={
    draw,
    rounded corners=2pt,
    align=left,
    minimum height=10mm,
    text width=37mm,
    inner sep=4pt,
    fill=white
  },
  widebox/.style={
    draw,
    rounded corners=2pt,
    align=left,
    minimum height=10mm,
    text width=50mm,
    inner sep=4pt,
    fill=white
  },
  datastore/.style={box, fill=blue!6, draw=blue!45!black},
  process/.style={box, fill=teal!6, draw=teal!50!black},
  memory/.style={box, fill=orange!10, draw=orange!60!black},
  infer/.style={box, fill=green!7, draw=green!45!black},
  output/.style={box, fill=purple!6, draw=purple!45!black},
  note/.style={box, fill=gray!10, draw=gray!50!black},
  line/.style={-Latex, thick},
  dashedline/.style={-Latex, thick, dashed}
]

% Left: dataset and sample construction
\node[datastore] (datasets) at (0, 0) {\textbf{Forecasting 数据源}\\
三类 benchmark 数据集\\
\quad 1. tourism\_monthly\\
\quad 2. nn5\_daily\\
\quad 3. australian\_electricity\_demand\\
文件格式：\texttt{.tsf}，包含频率、horizon、属性与序列值};

\node[process, below=of datasets] (builder) {\textbf{样本构造：\texttt{build\_tsf\_forecasting\_dataset}}\\
读取 \texttt{.tsf} 元数据与序列\\
窗口切分、历史长度设定、尺度归一化\\
生成 train / val / test 样本};

\node[widebox, below=of builder, fill=blue!3, draw=blue!40!black] (sample) {\textbf{\texttt{ForecastSample}}\\
\texttt{sequence}（归一化窗口）\\
\texttt{static}（均值、尺度、last、trend、seasonal gap、history length、horizon）\\
\texttt{raw\_context / raw\_target}（参数无关经验写入基础）\\
\texttt{metadata}（series, seasonality, window\_end\_index 等）};

\node[process, right=28mm of builder] (formation) {\textbf{Formation 抽取：\texttt{build\_window\_formation}}\\
16 个 formation 特征\\
\quad 趋势、波动、季节、level shift、相位、曲率、自相关、稳定性等\\
弱标签生成：\\
\texttt{pattern\_label}\\
\texttt{trajectory\_label}\\
\texttt{experience\_label}};

% Encoder
\node[process, right=28mm of sample, yshift=8mm] (encoder) {\textbf{共享编码主干：GRU Manifold Encoder}\\
输入：\texttt{sequence} + \texttt{static}\\
时序骨干：双向 GRU\\
多视角上下文：\\
\quad \texttt{final\_hidden}\\
\quad \texttt{pooled\_hidden}\\
\quad \texttt{delta\_hidden}\\
\quad \texttt{static\_projector(static)}};

\node[note, below=of encoder] (context) {\textbf{Context 表示}\\
\[
\texttt{context} = [h_{\text{final}}; h_{\text{pool}}; h_{\Delta}; h_{\text{static}}]
\]
不是单一最后隐藏态，而是末端状态、全局轮廓、变化方向与静态背景的联合表示};

\node[process, below=of context] (qkv) {\textbf{投影头：\texttt{query / key / value / input\_embedding}}\\
\texttt{query}, \texttt{key}：L2 normalize，用于相似性检索\\
\texttt{value}：保持内容表达\\
\texttt{input\_embedding}：当前窗口主表示};

% Training time memory build
\node[process, right=32mm of encoder, yshift=28mm] (managerwrite) {\textbf{训练期记忆写入：\texttt{MetaMemoryManager.write}}\\
输入：编码结果 + 三个弱标签 + metadata\\
将样本分发到三类 memory 组件\\
并记录 activity / support / write\_confidence / freshness};

\node[memory, below=of managerwrite] (pmemory) {\textbf{PatternMemory}\\
局部形态原型\\
同 label 合并，偏波形相似};

\node[memory, below=of pmemory] (tmemory) {\textbf{TrajectoryMemory}\\
阶段走势原型\\
阈值更严，偏 regime 相似};

\node[memory, below=of tmemory] (ememory) {\textbf{ExperienceMemory}\\
综合经验原型\\
支持 hot bank / archive bank\\
支持 blended retrieval};

% Persistent store
\node[datastore, right=34mm of managerwrite, yshift=5mm] (store) {\textbf{Persistent Experience Store}\\
\texttt{upsert\_samples(...)}\\
保存参数无关经验\\
\texttt{load\_samples(...)} / \texttt{load\_prototypes(...)}\\
跨运行、跨实验经验复用};

\node[datastore, below=of store] (files) {\textbf{持久化文件层}\\
\texttt{experience\_entries.jsonl}\\
\texttt{prototypes.jsonl}\\
\texttt{events.jsonl}\\
\texttt{semantic\_index.json}\\
\texttt{cache/<model\_fingerprint>.pt}};

% Inference path
\node[infer, below=45mm of qkv] (readmanager) {\textbf{推理期调度：\texttt{MetaMemoryManager.read}}\\
场景特征：horizon ratio, seasonality ratio, series density, regime mix\\
planner weights：pattern / trajectory / experience\\
label priors：pattern, trajectory, experience\\
semantic retrieval：prototype top-k 命中};

\node[infer, right=28mm of readmanager] (attnread) {\textbf{三路读取：Attention Memory Reader}\\
\[
\texttt{query} \cdot \texttt{keys} \rightarrow \text{top-k} \rightarrow \text{softmax}(\texttt{values})
\]
输出：readout、matched labels、confidence、top label\\
experience 分支可融合 archive 与 prototype template};

\node[infer, right=28mm of attnread] (fusion) {\textbf{融合层：Gate Network + Multi Fusion}\\
门控输入：\\
\quad \texttt{input\_embedding}\\
\quad normalized static\\
\quad normalized formation\\
\quad planner weights\\
\quad memory confidence\\
\quad scenario features\\
输出：三路 gate 与融合表示};

\node[infer, below=of fusion] (heads) {\textbf{预测头}\\
\texttt{base\_regressor}：直接预测整段 horizon\\
\texttt{bucket\_heads}：short / mid / long 三段 residual\\
\[
\hat{y}_{fusion} = \hat{y}_{base} + \hat{y}_{residual}
\]};

\node[output, right=28mm of heads] (outputs) {\textbf{输出与诊断}\\
最终预测：\texttt{test\_metrics}\\
基线对照：\texttt{base\_only}\\
记忆增益：\texttt{memory\_effectiveness}\\
过程诊断：\texttt{memory\_diagnostics}\\
样例检索轨迹：\texttt{example\_retrieval\_trace}};

% Semantic prototype branch
\node[infer, above=of attnread] (prototype) {\textbf{Prototype 语义模板}\\
\texttt{semantic\_retrieve(...)}\\
formation similarity + label match + support + quality\\
聚合 \texttt{template\_curve}\\
经 \texttt{prototype\_curve\_encoder} 变成 semantic value};

% Neural cache branch
\node[datastore, below=of files] (cache) {\textbf{Neural Cache}\\
\texttt{compute\_model\_fingerprint()}\\
\texttt{export\_neural\_cache()}\\
\texttt{load\_neural\_cache(...)}\\
同模型条件下直接恢复最终 memory bank};

% Legend
\node[note, below=40mm of readmanager, text width=196mm] (legend) {\textbf{图注与学术说明}\\
\textbf{训练阶段：} 数据集窗口化后，经 formation 与 GRU manifold encoder 生成 \texttt{query/key/value/input\_embedding}，样本依据三个弱标签写入三类 memory，同时原始窗口与语义摘要写入 persistent experience store。\\
\textbf{推理阶段：} 当前窗口先由 manager 生成 planner weights 与 label priors，再联合 semantic prototype 与三路 memory 做检索，随后通过 gate network 和 residual forecasting heads 生成最终预测。\\
\textbf{系统特征：} 该架构同时包含参数无关经验层（entries / prototypes）、参数相关神经缓存层（fingerprint / cache）、以及 end-to-end 训练的 GRU retrieval-fusion forecasting 主干。};

% arrows left path
\draw[line] (datasets) -- (builder);
\draw[line] (builder) -- (sample);
\draw[line] (builder) -- (formation);
\draw[line] (sample.east) -- ++(8mm,0) |- (encoder.west);
\draw[line] (formation.south) |- (encoder.north);
\draw[line] (encoder) -- (context);
\draw[line] (context) -- (qkv);

% training arrows
\draw[line] ($(qkv.east)+(0,18mm)$) -- (managerwrite.west);
\draw[line] (formation.east) -- ++(10mm,0) |- (managerwrite.west);
\draw[line] (managerwrite) -- (pmemory);
\draw[line] (managerwrite) -- (tmemory);
\draw[line] (managerwrite) -- (ememory);
\draw[line] (sample.east) -- ++(20mm,0) |- (store.west);
\draw[line] (store) -- (files);
\draw[dashedline] (managerwrite.east) -- (store.west);
\draw[dashedline] (ememory.east) -- ++(10mm,0) |- (cache.west);
\draw[dashedline] (store.south) -- (cache.north);

% inference arrows
\draw[line] (qkv.south) -- ++(0,-10mm) -| (readmanager.north);
\draw[line] (formation.south) |- (readmanager.west);
\draw[line] (pmemory.east) |- (attnread.north west);
\draw[line] (tmemory.east) -- ++(12mm,0) |- (attnread.north);
\draw[line] (ememory.east) |- (attnread.south west);
\draw[line] (readmanager) -- (attnread);
\draw[line] (readmanager.east) -- ++(12mm,0) |- (prototype.west);
\draw[dashedline] (files.west) -- ++(-10mm,0) |- (prototype.north);
\draw[dashedline] (cache.west) -- ++(-14mm,0) |- (readmanager.south);
\draw[line] (prototype) -- (attnread);
\draw[line] (attnread) -- (fusion);
\draw[line] (fusion) -- (heads);
\draw[line] (heads) -- (outputs);

% group boxes
\begin{scope}[on background layer]
  \node[draw=black!45, rounded corners=4pt, dashed, inner sep=6pt, fit=(datasets)(builder)(sample)(formation), label={[font=\bfseries]north:阶段 A：数据准备与语义构造}] {};
  \node[draw=black!45, rounded corners=4pt, dashed, inner sep=6pt, fit=(encoder)(context)(qkv)(managerwrite)(pmemory)(tmemory)(ememory), label={[font=\bfseries]north:阶段 B：训练期表示学习与记忆形成}] {};
  \node[draw=black!45, rounded corners=4pt, dashed, inner sep=6pt, fit=(store)(files)(cache), label={[font=\bfseries]north:阶段 C：持久化经验层与神经缓存层}] {};
  \node[draw=black!45, rounded corners=4pt, dashed, inner sep=6pt, fit=(readmanager)(prototype)(attnread)(fusion)(heads)(outputs), label={[font=\bfseries]north:阶段 D：推理期检索、融合与预测}] {};
  \node[draw=black!45, rounded corners=4pt, dashed, inner sep=6pt, fit=(legend), label={[font=\bfseries]north:图例说明}] {};
\end{scope}

\end{tikzpicture}
\end{document}
